\section{Files}

\begin{frame}[fragile]{File reading}
  Put input data in a text / YAML file (modern approach in this repo):
\begin{lstlisting}[style=terminal]
# input.yaml
v0: 2.0
a: 0.2
dt: 0.1
interval: [0, 2]
\end{lstlisting}
  \begin{alertblock}{Question}
    How can we read this file into variables \texttt{v0}, \texttt{a},
    \texttt{dt}, and \texttt{interval}?
  \end{alertblock}
\end{frame}

\begin{frame}[fragile]{Reading parameters with YAML}
\lstinputlisting[style=pythonTiny,firstline=7,lastline=26]{code/for_file_read_input.py}
\end{frame}

\begin{frame}[fragile]{Classic line parsing: \texttt{variable = value}}
  Older style (still useful to understand):
\begin{lstlisting}[style=pythonTiny]
with open('input.dat') as infile:
    for line in infile:
        name, value = line.split('=')
        name = name.strip()
        if name in ('v0', 'a', 'dt'):
            exec(f'{name} = float({value})')
        elif name == 'interval':
            interval = eval(value)
\end{lstlisting}
\end{frame}

\begin{frame}[fragile]{Splitting lines into words}
\begin{lstlisting}[style=pythonTiny]
>>> line = 'v0 = 5.3'
>>> variable, value = line.split('=')
>>> variable
'v0 '
>>> value
' 5.3'
>>> variable.strip()  # strip away blanks
'v0'
\end{lstlisting}
  Note: must convert \texttt{value} to \texttt{float} before computing!
\end{frame}

\begin{frame}[fragile]{The magic \texttt{eval} function}
  \texttt{eval(s)} executes a string \texttt{s} as a Python expression
  and creates the corresponding object:
\begin{lstlisting}[style=pythonTiny]
>>> obj1 = eval('1+2')
>>> obj1, type(obj1)
(3, <class 'int'>)
>>> obj2 = eval('[-1, 8, 10, 11]')
>>> obj2, type(obj2)
([-1, 8, 10, 11], <class 'list'>)
>>> from math import sin, pi
>>> x = 1
>>> obj3 = eval('sin(pi*x)')
>>> obj3, type(obj3)
(1.2246467991473532e-16, <class 'float'>)
\end{lstlisting}
\end{frame}

\begin{frame}{Why is \texttt{eval} useful?}
  We can read formulas, lists, expressions as text from a file and with
  \texttt{eval} turn them into live Python objects.

  \begin{alertblock}{Caution}
    Only use \texttt{eval} on trusted input (never on arbitrary user
    strings from the web).
  \end{alertblock}
\end{frame}

\begin{frame}[fragile]{Implementing a calculator in Python}
  Demo:
\begin{lstlisting}[style=terminal]
Terminal> uv run python calc.py "1 + 0.5*2"
2.0
Terminal> uv run python calc.py "sin(pi*2.5) + exp(-4)"
1.0183156388887342
\end{lstlisting}
  Just a few lines:
\begin{lstlisting}[style=pythoncompact]
import sys
from math import *   # sin, cos, exp, pi, etc.
expression = sys.argv[1]
result = eval(expression)
print(result)
\end{lstlisting}
\end{frame}

\begin{frame}[fragile]{The \texttt{with} statement for file handling}
\begin{lstlisting}[style=pythoncompact]
with open('input.dat', 'r') as infile:
    for line in infile:
        ...
\end{lstlisting}
  No need to close the file when using \texttt{with}.
\end{frame}

\begin{frame}[fragile]{File writing}
  \begin{itemize}
    \item We have $t$ and $s(t)$ in lists \texttt{t\_values}, \texttt{s\_values}
    \item Task: write a nicely formatted table to a file
  \end{itemize}
\begin{lstlisting}[style=pythonTiny]
with open('table1.dat', 'w') as outfile:
    outfile.write('# t    s(t)\n')
    for t, s in zip(t_values, s_values):
        outfile.write(f'{t:.2f}  {s:.4f}\n')
\end{lstlisting}
\end{frame}

\begin{frame}[fragile]{Simplified writing: \texttt{numpy.savetxt}}
\begin{lstlisting}[style=pythonTiny]
data = np.array([t_values, s_values]).transpose()
np.savetxt('table2.dat', data, fmt=['%.2f', '%.4f'],
           header='t   s(t)', comments='# ')
\end{lstlisting}
  Resulting \texttt{table2.dat}:
\begin{lstlisting}[style=terminal]
# t   s(t)
0.00 0.0000
0.10 0.2010
0.20 0.4040
...
2.00 4.4000
\end{lstlisting}
\end{frame}

\begin{frame}[fragile]{Simplified reading: \texttt{numpy.loadtxt}}
\begin{lstlisting}[style=pythoncompact]
data = np.loadtxt('table2.dat', comments='#')
\end{lstlisting}
  Note:
  \begin{itemize}
    \item Lines beginning with \texttt{\#} are skipped
    \item \texttt{data} is a 2-D array: \texttt{data[i,0]} holds $t$
          and \texttt{data[i,1]} holds $s(t)$ in row $i$
  \end{itemize}
  Full script: \texttt{python\_intro/for\_file\_read\_input.py}
\end{frame}
